\clearpage
\appendix
\section{Author notes}
\textbf{Acknowledgements.} I thank Professor Luyao Zhang. The PS1 final review asked for a two-sided game and visible matrices; this final revision implements both. On 2026-09-29 I withdrew the chatbot-subsidy-auction draft because it did not continue the PS1. The 28 September discussion was about a withdrawn draft of a different model. Those comments are not reviews of this paper, so they are not quoted here.

\textbf{Contribution.} Muhan Chen is the sole author. My intellectual decisions were to make the Promiser a real player, set the forfeit at $5$ because the gain from Break is $4$, keep similarity outside the matrix, and make two humans play the game. Outputs are the GitHub, Colab, and Hugging Face artifacts. Verification is the dominance proof, the three Nashpy checks, and the seed-206 replay. The human sample is seven pairs in the repository folder \texttt{data/}.

\bibliographystyle{ACM-Reference-Format}
\bibliography{references}

\section{Parameters and reproducibility}
\textbf{Players.} Two humans: Promiser (Keep/Break) and Decider (Trust/Verify). \textbf{Score cards.} Aligned: Promiser Keep/Break payoffs $6/3$; misaligned: $6/10$ before the forfeit. \textbf{Labels.} SAME if the displayed priorities match; DIFFERENT otherwise. The label changes no payoff and cannot waive the forfeit. \textbf{Decider payoffs.} Trust/Keep $6$, Trust/Break $-4$, Verify $2$. \textbf{Promiser payoffs.} Aligned: Keep $6$, Break $3$; misaligned: Keep $6$, Break $10$, or $5$ after the forfeit. \textbf{Information.} Public card and label; actions hidden until both lock; simultaneous action. \textbf{Horizon.} Twelve fresh rounds; no carry-over score. \textbf{Benchmark.} $10p-4=2$, so $p^\ast=0.6$. \textbf{Computation.} Nashpy for matrices; Python/Colab for replay; interface test marked simulated.

Keep is strictly better than Break on the aligned card ($6>3$). On the misaligned card without a forfeit, Break is strictly better ($10>6$), so the Decider's best response is Verify ($2>-4$). With the $5$ forfeit, Keep is strictly better ($6>5$), and Trust is then the Decider's best response ($6>2$). Nashpy returns exactly the three pure equilibria named in Section 3. Reproduction uses commit \texttt{b1380f7}, the matrix test, and Colab Runtime $\rightarrow$ Run all; the saved seed-206 output matches 48/24/22/18. The full hash is in the links section.

\textbf{AI-use disclosure.} 2026-09-29 to 2026-09-30, Cursor coding agent. Purpose: repository, notebook, Hugging Face game, and revision prompts from the author's locked model. Accepted: human Promiser and Decider, forfeit of 5, the three matrices, and Nashpy. Rejected: an AI player, a personal Drive Colab citation, and invented human data.

\section{Pseudocode and horizon record}
\textbf{Pseudocode.} (1) Read the score card and whether the forfeit rule is active. (2) If the card is aligned, set the forfeit to false; SAME/DIFFERENT never changes it. (3) Fill the $2\times2$ payoff matrices from the stated Decider and Promiser payoffs. (4) Use strict dominance and best responses to mark the unique pure Nash equilibrium, or record ``not unique'' at a $4$-point second bid. (5) Replay seed 206 and record PS1 totals 48, 24, 22, and 18. (6) For the appendix comparison, record only the second bid and return the $>4$, $<4$, or $=4$ outcome.

\textbf{PS1.} $p^\ast=0.6$; Keep was a draw. \textbf{PS2.} Forfeit 5 flips the matrix. In seven pairs, Trust was higher on forfeit cards, descriptively and mixed with order; Keep did not move. \textbf{Next.} Do not treat interface clicks as behavior. \textbf{2056.} A checked rule for when a copied style cannot replace the payoff card.

The 2056 hope is to be remembered for a testable rule: in public or commercial negotiation, similar style cannot replace the payoff card's forfeit. It builds on the PS1 threshold, these three equilibria, and seven exploratory pairs, not on a claim of a Nobel prize or completed SDG impact.

\section{Four-source literature record}
\textbf{Chen, PS1, 2026.} $p^\ast=0.6$; similarity changes no payoff; Keep is a draw. This paper lets the Promiser choose, and a forfeit of 5 flips the misaligned equilibrium. \textbf{Crawford and Sobel 1982.} Cheap talk need not change the Break payoff. The forfeit changes that cell, and the label cannot waive it. \textbf{Berg, Dickhaut, and McCabe 1995}~\cite{berg1995}. The second mover chooses a return. Here the card makes Keep or Break strictly dominant before play. \textbf{Nash 1950; Osborne and Rubinstein 1994.} The equilibrium definition is applied to these three matrices and checked in Nashpy.

\section{PS1 baseline and revision}
PS1's threshold was $p^\ast=0.6$ and its similarity cue did not alter payoffs. Its Keep probabilities were exogenous. The final review dated 2026-09-25 reports 75.57/100, Research Quality $Q=59/100$, Research $50.57/60$, and Revision $0/15$ because an implemented v1-to-v2 change was not visible. The seed-206 totals copied from the graded PS1 are benchmark $+48$, always-Verify $+24$, always-Trust $+22$, and Trust-if-similar $+18$. The present revision makes Promiser action endogenous and adds the $5$ misaligned-card forfeit; the result changes from a decision threshold to a Nash prediction.

\section{Pair-level human record}
Each row is one file in \texttt{data/}. Trust and Keep are counts out of 12 rounds. The last column is Keep on the four forfeit cards minus Keep on the four misaligned cards with no forfeit.

Pair 1: Trust 9/12, Keep 6/12, forfeit minus no forfeit $-1$. Pair 2: 7/12, 8/12, $-1$. Pair 3: 6/12, 7/12, $0$. Pair 4: 9/12, 8/12, $+2$. Pair 5: 8/12, 5/12, $0$. Pair 6: 10/12, 4/12, $+1$. Pair 7: 11/12, 6/12, $-1$.

\section{Three-lens stress test}
\textbf{Game theory.} If many Promisers Keep without the forfeit, strict dominance does not describe them. That check is triggered: on misaligned cards with no forfeit they Keep in 15/28 rounds, and with the forfeit of 5 they still Break in 13/28. \textbf{Social choice.} If the Decider's session score does not reach the equilibrium total of 56, that trade-off is not what these pairs received. Their mean session score is 19.7. \textbf{Mechanism design.} The label cannot waive the forfeit in the rules. They Kept in only 15/28 rounds when the forfeit was on the card. The rule was displayed; the Promiser's behavior did not implement it.

\section{Revision map from the final PS1 review}
\textbf{Matrix.} Three matrices are printed in Section 2, including the forfeit-of-5 card. The README prints the same three. Commit \texttt{b1380f7}. \textbf{Both sides act.} The Promiser now chooses Keep or Break. Seven pairs are in Section 4. \textbf{Information.} The card is public and the actions are hidden until both lock. Scores do not carry over. \textbf{Nash.} Strict dominance, best responses, and deviations are in Sections 2--3. Only pure strategies are claimed. \textbf{Nashpy.} The output matches all three matrices. \textbf{Separate evidence.} Section 4 separates the equilibrium replay, interface clicks, and seven human pairs. The 84 rounds are valid; the intervals are descriptive. \textbf{One commit.} The course repository, Colab notebook, and Hugging Face Space use the hash in the links section. \textbf{Keep.} The two human roles and the matrices stay. The seven pairs do not replace the matrix. \textbf{Revise.} Figure 1 is the evidence chain. Trust was higher on forfeit cards, descriptively and mixed with order; Keep did not move. \textbf{Investigate.} An Innovation Studio run was not completed, so no Studio URL is cited. The benchmark, forfeit cell, overturning condition, and illustrative-score limit are in the evidence section.

\section{Evidence and withdrawn draft}
The benchmark is $p^\ast=0.6$, and Trust follows only when Keep is strictly better. Cheap talk has no forfeit cell, and PS1's Promiser had no choice. The flip fails if Break's gross payoff is greater than $11$ or the forfeit cannot be collected. The payoff numbers in the matrices are illustrative. The human scores are observed. The sample is seven pairs. The withdrawn chatbot-subsidy-auction draft is not the model of this paper.

\section{Links and open-science checklist}
The final release uses the same organization commit in \GitHubLink, \NotebookLink, and \HuggingFaceLink. The full commit is \texttt{b1380f743b278819d514ec24a36aef89eb062829}. The repository contains the matching matrices, code, parameters, tests, outputs, and MIT license. The notebook exposes the seed, runs the matrix checks, and reproduces $+48/+24/+22/+18$. The Space is the two-player interface. The seven session files are the human sample. Other clicks are not that sample.

\section{Eight-row structure summary}
\textbf{Background.} This paper: two humans choose Keep/Break and Trust/Verify around an aligned or misaligned card plus a separate similarity label. PS1: Trust/Verify around a public payoff table and a similarity cue. \textbf{Motivation.} This paper: a style-matching negotiation screen can invite trust when the card is misaligned; a forfeit tests a correction. PS1: trust may follow resemblance rather than incentives. \textbf{Questions.} This paper: the two misaligned equilibria, whether computation reproduces them and PS1, and whether people follow the card or SAME. PS1: when the trust threshold applies, and whether similarity changes behavior. \textbf{Method.} This paper: simultaneous hidden actions, public payoffs, twelve fresh rounds, three matrices, Nashpy, and seven sessions. PS1: exogenous Keep probabilities $0.85/0.25$ and a seed-206 replay. \textbf{Results.} Actual: three Nash equilibria; PS1 totals $+48/+24/+22/+18$; seven pairs, Decider match 52/84, Promiser match 42/84. PS1 had no human data. \textbf{Merits.} A $5$ loss changes strict dominance and leaves similarity outside the table. PS1 separates incentive-based trust from similarity-based over-trust. \textbf{Impacts.} Disclosure and a non-waivable forfeit are questions, including SDG 10 and 16. PS1 impacts were not measured. \textbf{Revision.} PS1 feedback led to a two-sided human game, visible matrices, Nashpy parity, corrected links, and a split between model output and human play. The review had identified missing strategic actions, a missing matrix, an unclear information structure, and missing model-aligned computation.

\section{Auction record}
This is a bounded allocation record in the appendix, not the paper's contribution. Two prospective Promisers each submit a nonnegative integer bid for how much they are willing to have taken if they Break on a misaligned card. The two bids compete only to determine who becomes Promiser. The Decider is not selected by the bid. The bid is a private willingness signal, while the score card and SAME/DIFFERENT label are public. SAME/DIFFERENT cannot change a bid or break a tie; a session seed breaks ties.

The higher bidder's bid determines the role, but the forfeit written into a misaligned Break cell is the second-highest bid. It is paid only after Break, not after Keep. If the second bid is greater than $4$, the Break payoff is $10$ minus that bid and is strictly below Keep's $6$, so the unique equilibrium is Trust, Keep. If the second bid is less than $4$, Break is strictly better than Keep and the unique equilibrium is Verify, Break. If the second bid equals $4$, both Keep and Break give the Promiser $6$, so there is no unique equilibrium. These cases follow the code's rule exactly; no claim is made that bidding cost is dominant or that this is a standard independent-private-values auction.

\textbf{Values or signals.} Each bid is a nonnegative integer. The card and the SAME/DIFFERENT label are public. \textbf{Prediction.} Second bid $>4$: Trust, Keep; $<4$: Verify, Break; $=4$: no unique equilibrium. \textbf{Bid and allocation.} The higher bidder becomes Promiser. The bid does not choose the Decider. A seed breaks ties. \textbf{Payment.} On misaligned Break, the Promiser pays the second-highest bid. On Keep, payment is zero. \textbf{Stopping.} One sealed pair of bids, then one hidden Keep/Break and Trust/Verify choice. \textbf{Utility.} Decider gets $6$, $-4$, or $2$. Promiser gets $6$ on Keep and $10$ minus the forfeit on Break. \textbf{Revenue.} If the second bid is greater than $4$, equilibrium is Keep and nothing is collected. If it is less than $4$, equilibrium is Break and revenue equals the second bid. \textbf{Efficiency.} Role assignment ignores similarity. Fairness is unresolved. \textbf{Behavior.} This auction was not played. The card game has seven pairs. If a later bid sample follows SAME, revise the protocol and keep the matrix.
